%% kapa: exam=1 round=27 year=2016 period=2 form=A subject=accounting subject_ko=회계학 seq=24 number=64
%% source: Q-Net 감정평가사 기출문제 (cst00309) artlSeq=5169402 file=extracted/1008/5169402/제27회 감정평가사 1차 문제/제27회 감정평가사 2교시문제(A형).pdf
%% conversion: kordoc 4.24.51 -> tools/gen1.py (auto); review=proofread-sonnet; flags=none
\begin{question}{64}
일반목적재무보고의 목적에 관한 설명으로 옳지 않은 것은?
\begin{choices}
  \item 현재 및 잠재적 투자자, 대여자 및 기타 채권자가 기업에 자원을 제공하는 것에 대한 의사결정을 할 때 유용한 보고기업 재무정보를 제공하는 것이다.
  \item 지분상품 및 채무상품을 매수, 매도 또는 보유하는 것에 대한 현재 및 잠재적 투자자의 의사결정은 그 금융상품 투자에서 그들이 기대하는 수익, 예를 들어, 배당, 원금 및 이자의 지급 또는 시장가격의 상승에 의존한다.
  \item 경영진의 책임 이행에 대한 정보는 경영진의 행동에 대해 의결권을 가지거나 다른 방법으로 영향력을 행사하는 현재 투자자, 대여자 및 기타 채권자의 의사결정에도 유용하다.
  \item 일반목적재무보고서는 보고기업의 가치를 보여주기 위해 고안된 것이다. 따라서 그 보고서는 현재 및 잠재적인 정보이용자가 보고기업의 가치를 추정하는데 도움이 되는 정보를 제공한다.
  \item 보고기업의 경영진도 해당 기업에 대한 재무정보에 관심이 있다. 그러나 경영진은 그들이 필요로 하는 재무정보를 내부에서 구할 수 있기 때문에 일반목적재무보고서에 의존할 필요가 없다.
\end{choices}
\end{question}
\answer{4}
